In his paper, Goldsmith describes a possible extension to bivariate data.  
He concatenates data and spline coefficient matrices.  
Specifically, let \(Y(t) = \left[Y_1(t),Y_2(t)\right]\) and let us express this in matrix form as \(Y = [Y_1,Y_2]\), 
where \(Y_i, i\in(1,2)\) corresponds to functional outcomes in the X and Y dimensions respectively concatenated together by 
column, \(B_W^T = [B_{1,W}^T, B_{2,W}^T]\), corresponds to the fixed effect spline coefficients in two dimensions, 
and \(B_Z^T = [B_{1,Z}^T, B_{2,Z}^T]\) corresponds to the random effect spline coefficients in two dimensions.  
Taking advantage of Kronecker product structure\footnote{See appendix},  Goldsmith writes out the following model:

\begin{equation}
    \begin{split}
        Y = ZB_Z^T\left(I_2 \otimes\Theta^T\right) + \epsilon\\
        \epsilon\sim N[0, \Sigma\otimes I_n]\\
        \Sigma \sim IW[v, \Psi]\\
        B_{Z,i} \sim N[B_w \boldsymbol{w}_i^T, \text{diag}\left(\sigma_{1,Z}^2, \sigma_{2,Z}^2\right)\otimes P^{-1}]\quad \text{for} \quad i = 1,\dots, I\\
        \sigma_{m,Z}^2 \sim IG[a_{m,Z}, b_{m,Z}]\quad \text{for} \quad m = 1,2\\
        B_{W,i} \sim N[0, \text{diag}\left(\sigma_{1,W_{k}}^2, \sigma_{2,W_{k}}^2\right)\otimes P^{-1}]\\
        \sigma_{m,W_{k}}^2 \sim IG[a_{m,W_{k]}}, b_{m,W_{k}}]\quad \text{for} \quad m = 1,2 \quad\text{and} \quad k = 1,\dots, p.
    \end{split}
\end{equation}

Note how fixed and random effects in each dimension are assumed independent.  
This bivariate model allows for modeling position curves in two dimensions concurrently.  
This bivariate model explicitly models correlation between the X and Y position curves through the level 1 residual 
covariance matrix; the X and Y dimensions are independent in both the fixed and random effects.      

Goldsmith remarks \say{Similar extensions to three or more curves \ldots proceed similarly}.  
We continue to follow Goldsmith's lead, giving the following model for three dimensional functional data:

\begin{equation}
    \begin{split}
        Y = ZB_Z^T\left(I_3 \otimes\Theta^T\right) + \epsilon\\
        \epsilon\sim N[0, \Sigma\otimes I_n]\\
        \Sigma \sim IW[v, \Psi]\\
        B_{Z,i} \sim N[B_w \boldsymbol{w}_i^T, \text{diag}\left(\sigma_{1,Z}^2, \sigma_{2,Z}^2, \sigma_{2,Z}^3\right)\otimes P^{-1}]\quad \text{for} \quad i = 1,\dots, I\\
        \sigma_{m,Z}^2 \sim IG[a_{m,Z}, b_{m,Z}]\quad \text{for} \quad m = 1,2, 3\\
        B_{W,i} \sim N[0, \text{diag}\left(\sigma_{1,W_{k}}^2, \sigma_{2,W_{k}}^2, \sigma_{3,W_{k}}^2\right)\otimes P^{-1}]\\
        \sigma_{m,W_{k}}^2 \sim IG[a_{m,W_{k]}}, b_{m,W_{k}}]\quad \text{for} \quad m = 1,2, 3 \quad\text{and} \quad k = 1,\dots, p
    \end{split}
\end{equation}

As with the bivariate model, fixed and random effects in each dimension are assumed to be independent.  
This independence assumption is important for two main reasons, one conceptual and another computational.  
First, assuming independence simplifes the covariance term in \(B_{Z,i}\) and \(B_{W,i}\).  
Without this independence assumption, we would need to model the correlation between X and Y fixed and random effect spline 
coefficients\footnote{We do this in the following model.}.  Secondly, assuming independence gives clear intution to the level 1 
residual covariance matrix:  it describes the manner in which the X and Y trajectories vary, after accounting for 
fixed and random effects\footnote{I think technically this would be the cross-variance.}. This trivariate model will serve 
as our first attempt at a trivariate model.  

\subsection{Concerns and Resulting Alternative Model}
Goldsmith's 2016 paper heavily motivates this work.  In his model (and its two dimensional extension), 
he places an inverse-Wishart prior on the level 1 covariance structure.  There is one cause for concern with this: 
for observations in three dimensions with each dimension being observed across a grid of \(D\), the level 1 covariance matrix 
would have to be \(3*D\times 3*D\).  In the case of \(D=100\), that would mean that our level 1 covariance matrix would be 
\(300\times300\).  If the level 1 covariance matrix is modeled as inverse Wishart, then as the dimension of this matrix increases, 
the degrees of freedom of the inverse Wishart also must increase for the distribution to remain proper. 
A result of this is that the Inverse Wishart \(IW\left(n\boldsymbol{S} + \Psi, n + v\right)\) posterior distribution tends to 
concentrate more tightly around its prior expected value: \(\Psi/{\left(v-p-1\right)}\)\cite{zhang_note_2021}. When the scale matrix is 
the identity matrix, this expected value is proportional to the identity matrix itself; using an Inverse Wishart prior with an identity 
scale matrix can lead to over-shrinkage of the estimated covariance matrix towards the identity matrix.